\subsubsection{Shao and Wu: Moment Contraction for Nonlinear Time Series}
\label{sec:shao-wu-gmc}

\paragraph{Published result.}
\citet{shao2007spectral} develop spectral theory for nonlinear time series.
Section~5, Theorem~5.1 (pp.~11--12 of the linked electronic reprint), provides
the relevant generation result. With i.i.d. innovations and a measurable,
time-homogeneous recurrence $\mathcal R_{\mathcal T}(\mathbf{s};\varepsilon)$,
its second-moment specialization assumes a finite second moment at some
reference history and the global bound
\begin{equation}
\begin{split}
 \|\mathcal R_{\mathcal T}(\mathbf{s};\varepsilon)
       -\mathcal R_{\mathcal T}(\mathbf{s}';\varepsilon)\|_{L^2}
 &\leq \sum_{d=0}^{L-1}a_d^{\mathrm{GMC}}|s_d-s'_d|,\\
 a_d^{\mathrm{GMC}}\geq0,\qquad
 \sum_{d=0}^{L-1}a_d^{\mathrm{GMC}}&<1.
\end{split}
\label{eq:gmc-generator}
\end{equation}
Here $\|Z\|_{L^2}=(\mathbb E[\|Z\|_2^2])^{1/2}$. The theorem yields a
causal strictly stationary solution and geometric moment contraction (GMC):
the expected squared difference between coupled versions with a shared future
innovation sequence decays geometrically. Vector-valued states are admitted;
the absolute differences become norms of lag blocks. Theorem~5.2 preserves
GMC under stable ARMA filtering, and Examples~5.1--5.5 cover exponential
autoregression, conditional heteroscedasticity, and bilinear recursions.

\paragraph{Benchmark adaptation.}
The standard factor mode wraps lag leaves involved in nonlinear products
with $\tanh$, retains standalone linear terms, and scales the AST until
its global sensitivity budget is below one. For example, the drift
$0.3\tanh(x_t)\tanh(x_{t-1})$ has
derivative bounds $0.3$ for each lag, giving a total bound of $0.6$; Gaussian
innovations satisfy the required moment condition. This route targets
stationary nonlinear dynamics on an unbounded state space. Global derivative
bounds can overestimate sensitivity when repeated factors represent the
same historical coordinate. The alternative standard mode applies $\tanh$
after a complete product and checks the uniform amplitude moment bound
derived in Section~\ref{sec:implemented-shao-wu}. That bound is a project
construction with its own stated conclusion. The paper's GMC conclusion
applies when its global sensitivity and moment hypotheses are verified.
A finite
direct uniform initialization can leave a finite-horizon transient; its law
is recorded, and no preliminary path is simulated or discarded. The AST
sampling and certification rules are project adaptations; deriving compatible
explanation labels from the retained equations is a further benchmark step.
